\documentclass[10pt,a4paper]{article}
\usepackage[utf8]{inputenc}
\usepackage[indonesian]{babel}
\usepackage{amsmath,amssymb,amsfonts}
\usepackage{geometry}
\usepackage{multicol}
\usepackage{enumitem}
\usepackage{fourier}

\geometry{top=1.5cm, bottom=1.5cm, left=1.5cm, right=1.5cm}

\pagestyle{empty}

\begin{document}

\begin{center}
  \textbf{\Large Latihan Soal Ulangan ASTS PABP}
\end{center}
\vspace{0.3cm}

\begin{multicols}{2}
  \setlist[enumerate,1]{leftmargin=*,label=\arabic*.}
  \setlist[enumerate,2]{leftmargin=*,label=\Alph*.}
  
  \begin{enumerate}
    \item Peristiwa yang menjadi \textit{asbabun nuzul} (sebab turunnya) QS. Ali Imran ayat 190-191 berkaitan dengan...
    \begin{enumerate}
      \item Pertanyaan kaum Yahudi mengenai ruh
      \item Permintaan kaum Quraisy kepada Nabi agar mengubah bukit Safa menjadi emas sebagai mukjizat
      \item Perintah hijrah ke Madinah
      \item Teguran kepada pasukan pemanah saat perang Uhud
      \item Larangan memakan riba berlipat ganda
    \end{enumerate}

    \item Kandungan dari penggalan lafal \textit{"\dots wakhtilāfil-laili wan-nahāri\dots"} dalam QS. Ali Imran ayat 190 adalah tentang...
    \begin{enumerate}
      \item Penciptaan bumi dan gunung
      \item Ciri-ciri orang yang berakal
      \item Pergantian malam dan siang
      \item Penciptaan langit bersap-sap
      \item Siksa api neraka
    \end{enumerate}

    \item Lanjutan hafalan lafal ayat tentang penciptaan alam semesta dalam QS. Ali Imran ayat 190 \textit{"Inna fī khalqis-samawāti wal-arḍi\dots"} adalah...
    \begin{enumerate}
      \item wakhtilāfil-laili wan-nahāri
      \item rabbanā mā khalaqta hāzā bāṭila
      \item la’āyātil li’ulil-albāb
      \item qiyāmaw wa qu‘ūdaw 
      \item fa qinā ‘azāban-nār
    \end{enumerate}

    \item Lanjutan hafalan lafal doa berlindung dari azab neraka pada QS. Ali Imran ayat 191 \textit{"\dots rabbanā mā khalaqta hāzā bāṭila, subḥānaka\dots"} adalah...
    \begin{enumerate}
      \item innal-‘ahda kāna mas’ūlā
      \item la’azīdannakum
      \item fatahlikū
      \item fa qinā ‘azāban-nār
      \item wa aufū bil-‘ahdi
    \end{enumerate}

    \item Hukum bacaan tajwid yang terdapat pada lafal \textit{"\dots s-samawāti"} dalam ayat \textit{"Inna fī khalqis-samawāti"} adalah...
    \begin{enumerate}
      \item Alif Lam Qamariyah
      \item Alif Lam Syamsiyah
      \item Izhar Halqi
      \item Idgham Bighunnah
      \item Ikhfa Syafawi
    \end{enumerate}

    \item Pada lafal \textit{"qiyāmaw wa"} dalam kalimat \textit{"Allazīna yazkurūnallāha qiyāmaw wa qu‘ūdaw"}, terdapat hukum bacaan...
    \begin{enumerate}
      \item Izhar Syafawi
      \item Idgham Bilaghunnah
      \item Iqlab
      \item Idgham Bighunnah
      \item Ikhfa Hakiki
    \end{enumerate}

    \item Pada huruf Nun berharakat dhammah di lafal \textit{"junūbihim"} (dalam penggalan \textit{wa qu‘ūdaw wa ‘alā junūbihim}), hukum bacaannya adalah...
    \begin{enumerate}
      \item Mad Thabi'i
      \item Mad Jaiz Munfasil
      \item Izhar Halqi
      \item Ikhfa Syafawi
      \item Iqlab
    \end{enumerate}

    \item Lafal dalam QS. Ali Imran ayat 191 yang menunjukkan ciri-ciri aktivitas \textit{ulil albab} (orang yang berakal) adalah...
    \begin{enumerate}
      \item Inna fī khalqis-samawāti wal-arḍi
      \item Tafakkarū fī khalqillāhi wa lā tafakkarū fī żātillāhi
      \item Allazīna yazkurūnallāha qiyāmaw wa qu‘ūdaw wa ‘alā junūbihim
      \item yatafakkarūna fī khalqis-samawāti wal-arḍ
      \item Yā ma‘syaral-jinni wal-insi inistaṭa‘tum
    \end{enumerate}

    \item Lengkapilah hafalan lafal hadits tentang larangan memikirkan Dzat Allah berikut: \textit{"Tafakkarū fī khalqillāhi wa lā \dots"}
    \begin{enumerate}
      \item yaqbiḍul-‘ilma biqabḍil-‘ulamā’i
      \item tafakkarū fī żātillāhi fatahlikū
      \item yasykurin-nāsa lā yasykurillāh
      \item yaqul khairan au liyaṣmut
      \item lā tanfużūna illā bisulṭān
    \end{enumerate}

    \item Bagaimana relasi antara iman, ilmu, dan amal berdasarkan kandungan QS. Ar-Rahman ayat 33?
    \begin{enumerate}
      \item Ilmu hanya wajib dituntut oleh para ulama
      \item Amal ibadah jauh lebih utama sehingga ilmu tidak diperlukan
      \item Jin lebih diutamakan dalam menuntut ilmu dibanding manusia
      \item Kecanggihan ilmu pengetahuan (IPTEK) harus semakin menumbuhkan keimanan kepada Allah
      \item Ilmu pengetahuan digunakan murni untuk mencari keuntungan materi
    \end{enumerate}

    \item Hukum tajwid pada lafal \textit{"wal-insi"} dalam QS. Ar-Rahman ayat 33 (huruf Nun sukun bertemu Sin) adalah...
    \begin{enumerate}
      \item Idgham Bighunnah
      \item Ikhfa Hakiki
      \item Iqlab
      \item Izhar Halqi
      \item Idgham Bilaghunnah
    \end{enumerate}

    \item Hukum tajwid pada akhir ayat lafal \textit{"\dots illā bisulṭān"} apabila diwaqafkan (diberhentikan) adalah...
    \begin{enumerate}
      \item Mad Arid Lissukun
      \item Mad Wajib Muttasil
      \item Mad Jaiz Munfasil
      \item Mad Iwad
      \item Mad Badal
    \end{enumerate}

    \item Lanjutan hafalan lafal hadits tentang dicabutnya ilmu \textit{"Innallāha lā yaqbiḍul-‘ilma-ntizā‘an yantazi‘uhū minal-‘ibādi, wa lākin \dots"} adalah...
    \begin{enumerate}
      \item falyaqul khairan au liyaṣmut
      \item fa qinā ‘azāban-nār
      \item innal-‘ahda kāna mas’ūlā
      \item yaqbiḍul-‘ilma biqabḍil-‘ulamā’i
      \item uḥillat lakum bahīmatul-an‘āmi
    \end{enumerate}

    \item Kandungan QS. Ar-Rahman ayat 33 menegaskan bahwa manusia dan jin tidak mampu menembus penjuru langit dan bumi kecuali dengan \textit{sulṭān}, yang bermakna...
    \begin{enumerate}
      \item Harta kekayaan yang melimpah
      \item Kekuatan militer yang besar
      \item Kekuatan / Penguasaan ilmu pengetahuan
      \item Izin dari para rasul
      \item Restu dari malaikat penjaga langit
    \end{enumerate}

    \item Contoh penerapan sikap berpikir kritis yang tepat bagi seorang pelajar di sekolah adalah...
    \begin{enumerate}
      \item Mengabaikan tugas kelompok yang diberikan guru
      \item Menerima semua informasi dari teman tanpa bertanya ulang
      \item Menghafal teori tanpa mau memahami maknanya
      \item Berani mempertanyakan materi pelajaran yang belum dipahami secara sopan
      \item Sering membolos karena merasa pelajaran terlalu sulit
    \end{enumerate}

    \item Cara menerapkan sikap berpikir kritis ketika menerima berita atau informasi dari media sosial adalah...
    \begin{enumerate}
      \item Langsung membagikan tautan berita tersebut ke grup lain
      \item Menyimpulkan isi berita hanya dari judulnya saja
      \item Mempercayai sepenuhnya karena berita tersebut sedang viral
      \item Menghapus dan memblokir pengirim tanpa alasan
      \item Melakukan klarifikasi (\textit{tabayyun}) sumber dan kebenaran isi berita
    \end{enumerate}

    \item Contoh penerapan berpikir kritis dalam mengamati fenomena alam di lingkungan sekitar adalah...
    \begin{enumerate}
      \item Meneliti penyebab terjadinya banjir bandang untuk mencari solusi pencegahan
      \item Membuang sampah di sungai karena merasa sudah biasa
      \item Menikmati keindahan alam sembari merusaknya perlahan
      \item Menganggap setiap bencana sebagai hal yang wajar tanpa perlu diantisipasi
      \item Acuh tak acuh terhadap perubahan cuaca di lingkungan rumah
    \end{enumerate}

    \item Bentuk penerapan sikap berpikir kritis pada saat harus mengambil sebuah keputusan penting adalah...
    \begin{enumerate}
      \item Mengikuti pendapat mayoritas agar aman, meskipun salah
      \item Memilih keputusan secara acak agar masalah cepat selesai
      \item Mempertimbangkan dampak baik dan buruk secara objektif sebelum memutuskan
      \item Menghindari tanggung jawab dari hasil keputusannya sendiri
      \item Meminta orang lain untuk memikirkan keputusannya
    \end{enumerate}

    \item Secara istilah, pengertian memenuhi janji adalah...
    \begin{enumerate}
      \item Mengucapkan kata-kata yang disenangi oleh banyak orang
      \item Melaksanakan sesuatu sesuai dengan kesanggupan yang telah diucapkan
      \item Memberikan hadiah kepada teman dekat secara sembunyi-sembunyi
      \item Menjaga amanah agar tidak diketahui oleh musuh
      \item Mengikuti semua keinginan orang tua tanpa membantah
    \end{enumerate}

    \item Lanjutan terjemahan hadits \textit{"Tanda-tanda orang munafik itu ada 3: jika berbicara ia berdusta, jika berjanji ia mengingkari, dan \dots"} adalah...
    \begin{enumerate}
      \item jika bersaksi ia palsu
      \item jika meminjam ia tidak mengembalikan
      \item jika berdagang ia curang
      \item jika diberi amanat ia berkhianat
      \item jika marah ia merusak
    \end{enumerate}

    \item Salah satu ciri utama orang yang selalu berusaha menepati janjinya adalah...
    \begin{enumerate}
      \item Mudah mengobral janji di setiap kesempatan
      \item Mencari-cari alasan agar terbebas dari tuntutan janji
      \item Menghindar saat ditagih janjinya oleh orang lain
      \item Berpikir terlebih dahulu mengenai kemampuannya sebelum berjanji
      \item Suka membatalkan janjinya secara sepihak
    \end{enumerate}

    \item Lanjutan hafalan lafal dalil Al-Quran tentang janji pada QS. Al-Isra ayat 34 \textit{"\dots wa aufū bil-‘ahdi, innal-‘ahda \dots"} adalah...
    \begin{enumerate}
      \item kāna mas’ūlā
      \item mā yurīd
      \item la’azīdannakum
      \item qaulan sadīdā
      \item raqībun ‘atīd
    \end{enumerate}

    \item Sikap yang seharusnya dilakukan apabila seseorang mengalami kendala di luar kendalinya sehingga tidak mampu memenuhi janji adalah...
    \begin{enumerate}
      \item Memutus kontak dengan orang yang telah diberi janji
      \item Berpura-pura lupa bahwa ia pernah membuat janji
      \item Bertanggung jawab dengan cara meminta maaf dan menjelaskan keadaannya
      \item Mencari kambing hitam atas kesalahannya tersebut
      \item Menunggu ditagih, lalu berpura-pura marah
    \end{enumerate}

    \item Contoh kasus nyata bentuk memenuhi janji seorang hamba kepada Allah SWT adalah...
    \begin{enumerate}
      \item Datang tepat waktu saat kerja bakti warga
      \item Membayar utang uang kepada pihak bank
      \item Melaksanakan salat 5 waktu dan menunaikan nazar berpuasa
      \item Mengembalikan buku perpustakaan sesuai batas peminjaman
      \item Mengerjakan tugas sekolah dengan teliti
    \end{enumerate}

    \item Contoh kasus perilaku menepati janji antar sesama manusia di lingkungan masyarakat adalah...
    \begin{enumerate}
      \item Melaksanakan ibadah kurban pada Hari Raya Idul Adha
      \item Membaca basmalah sebelum mulai bekerja
      \item Bersedekah di kotak amal masjid setiap salat Jumat
      \item Membayar barang pinjaman atau utang tepat waktu sesuai kesepakatan
      \item Menjauhi sikap iri dan dengki terhadap tetangga
    \end{enumerate}

    \item Definisi dari mensyukuri nikmat secara istilah adalah...
    \begin{enumerate}
      \item Merasa puas dengan apa yang ada dan berhenti berusaha
      \item Berbangga diri atas kekayaan yang dimiliki dan memamerkannya
      \item Menganggap kesuksesan yang didapat murni dari kerja keras diri sendiri
      \item Sikap menyadari bahwa segala nikmat berasal dari Allah dan menggunakannya sesuai kehendak-Nya
      \item Berbagi kekayaan untuk mendapat pujian masyarakat sekitar
    \end{enumerate}

    \item Lanjutan hafalan lafal dalil Al-Quran tentang bersyukur pada QS. Ibrahim ayat 7 \textit{"Wa iź ta’ażżana rabbukum la’in syakartum \dots"} adalah...
    \begin{enumerate}
      \item wa la’in kafartum
      \item inna ‘azābī lasyadīd
      \item kāna mas’ūlā
      \item lā yasykurillāh
      \item la’azīdannakum
    \end{enumerate}

    \item Penerapan sikap mensyukuri nikmat kesehatan yang telah diberikan oleh Allah SWT dilakukan dengan cara...
    \begin{enumerate}
      \item Berolahraga berat dan ekstrem tanpa memperdulikan kelelahan
      \item Bekerja keras siang dan malam tanpa waktu istirahat
      \item Menjaga pola makan, berolahraga, dan menggunakan tenaga untuk ibadah
      \item Mengonsumsi segala jenis makanan tanpa batasan asalkan enak
      \item Membiarkan tubuh kurang tidur demi bermain \textit{game}
    \end{enumerate}

    \item Cara dan penerapan sikap yang benar dalam mensyukuri nikmat harta benda adalah...
    \begin{enumerate}
      \item Menghamburkan uang untuk membeli berbagai koleksi barang mewah
      \item Menyimpan seluruh hartanya sendiri dan tidak mau bersedekah (bakhil)
      \item Mencari harta dengan segala cara, termasuk cara yang tidak halal
      \item Membelanjakan harta secukupnya, bersedekah, dan membantu yang membutuhkan
      \item Merusak barang yang masih bagus hanya karena bosan
    \end{enumerate}

    \item Secara istilah, definisi dari sikap menjaga lisan adalah...
    \begin{enumerate}
      \item Berhati-hati dalam berbicara dengan memilih ucapan yang baik, benar, bermanfaat, dan tidak menyakiti
      \item Memilih diam (tidak berbicara sama sekali) sepanjang hari
      \item Berani membicarakan keburukan orang lain secara terang-terangan
      \item Membela teman dekat dalam perdebatan meskipun teman tersebut bersalah
      \item Berkata jujur dengan bahasa yang sarkas dan menyinggung perasaan
    \end{enumerate}
  \end{enumerate}
\end{multicols}
\end{document}
